%%
%% This is file `pst-3dplot.tex',
%%
%% IMPORTANT NOTICE:
%%
%% Package `pst-3dplot.tex'
%%
%% Herbert Voss <voss@perce.de>
%%
%% This program can be redistributed and/or modified under the terms
%% of the LaTeX Project Public License Distributed from CTAN archives
%% in directory macros/latex/base/lppl.txt.
%%
%% DESCRIPTION:
%%   `pst-3dplot' is a PSTricks package to draw 3d curves and graphical objects
%%
\def\fileversion{1.30}
\def\filedate{2004/01/18}
\message{`PST-3dplot' v\fileversion, \filedate\space (Herbert Voss)}
\csname PSTThreeDplotLoaded\endcsname
\let\PSTThreeDplotLoaded\endinput
% Requires PSTricks, pst-node, pst-plot, multido packages
%\usepackage{pstcol}
\ifx\PSTricksLoaded\endinput\else\input pstricks.tex\fi
\ifx\PSTnodesLoaded\endinput\else\input pst-node.tex\fi
\ifx\PSTplotLoaded\endinput\else\input pst-plot.tex\fi% plotpoints
\ifx\PSTNewThreeDLoaded\endinput\else\input pst-vue3d.tex\fi
\ifx\PSTMultidoLoaded\endinput\else\input multido.tex\fi
%
% DPC interface to the `keyval' package (until keyval based version of PSTricks)
\input pst-key
%
\edef\PstAtCode{\the\catcode`\@} \catcode`\@=11\relax
\SpecialCoor
%
%%%%%%%%%%%%%%%%%% Macrolist %%%%%%%%%%%%%%%%%%%%%%%%%%%%%%%%%%%%
%
% \def\pstThreeDPut
% \def\pstThreeDLine
% \def\pstThreeDDot
% \def\pstThreeDNode
% \def\pstUThreeDPut
% \def\pstThreeDSquare
% \def\pstThreeDBox
% \def\pstThreeDSphere
% \def\pstThreeDEllipse
% \def\pstThreeDPlotFunc
% \def\psplotThreeD
% \def\parametricplotThreeD
% \def\psThreeDPlot
% \def\pstThreeDCoor
% \def\fileplotThreeD
% \def\dataplotThreeD
% \def\listplotThreeD
%
%
\newdimen\pst@dimf
%%%%%%%%%%%%%%%%%%%%%%%%%%%%%%%%%%%%%%%%%%%%%%%%%%%%%%%%%%%%
%
% ---------------- the if's -----------------
%
\newif\ifPst@drawing%			draw the coordinates?
\define@key{psset}{drawing}[true]{%
	\@nameuse{Pst@drawing#1}%
}
\newif\ifPst@drawCoor%		draw the coordinates of a dot?
\define@key{psset}{drawCoor}[false]{%
	\@nameuse{Pst@drawCoor#1}%
}
\newif\ifPst@hiddenLine%	emulate hidden line surface?
\define@key{psset}{hiddenLine}[false]{%
	\@nameuse{Pst@hiddenLine#1}%
}
\newif\ifPst@SphericalCoor% 	(r,phi,theta)
\define@key{psset}{SphericalCoor}[true]{\@nameuse{Pst@SphericalCoor#1}}
\setkeys{psset}{SphericalCoor=false}
%
% ------- the global definitions for the pspicture frame ------
%
\define@key{psset}{xMin}{\edef\psk@ThreeDplot@xMin{#1}}
\define@key{psset}{xMax}{\edef\psk@ThreeDplot@xMax{#1}}
\define@key{psset}{yMin}{\edef\psk@ThreeDplot@yMin{#1}}
\define@key{psset}{yMax}{\edef\psk@ThreeDplot@yMax{#1}}
\define@key{psset}{zMin}{\edef\psk@ThreeDplot@zMin{#1}}
\define@key{psset}{zMax}{\edef\psk@ThreeDplot@zMax{#1}}
\define@key{psset}{xThreeDunit}{\edef\psk@ThreeDplot@xThreeDunit{#1}}
\define@key{psset}{yThreeDunit}{\edef\psk@ThreeDplot@yThreeDunit{#1}}
\define@key{psset}{zThreeDunit}{\edef\psk@ThreeDplot@zThreeDunit{#1}}
%
% -------------- the angles and the plotpoints -------------
%
\define@key{psset}{Alpha}{%			Horizontal turn
	\edef\psk@ThreeDplot@Alpha{#1}%
}
\define@key{psset}{Beta}{%			Vertical turn
	\edef\psk@ThreeDplot@Beta{#1}%
}
\define@key{psset}{drawStyle}{%			how to draw 3D functions
	\edef\psk@ThreeDplot@drawStyle{#1}%
}
\def\drawStyle@xLines{xLines}
\def\drawStyle@yLines{yLines}
\def\drawStyle@xyLines{xyLines}
\def\drawStyle@yxLines{yxLines}
%
\define@key{psset}{xPlotpoints}{%		points in x direction
	\edef\psk@ThreeDplot@xPlotpoints{#1}%
}
\define@key{psset}{yPlotpoints}{%		points in y direction
	\edef\psk@ThreeDplot@yPlotpoints{#1}%
}
\define@key{psset}{beginAngle}{%		for ellipse/circle arc
	\edef\psk@ThreeDplot@beginAngle{#1}%
}
\define@key{psset}{endAngle}{%			for ellipse/circle arc
	\edef\psk@ThreeDplot@endAngle{#1}%
}
\define@key{psset}{linejoin}{%		how lines come together 0,1,2
	\edef\psk@ThreeDplot@linejoin{#1}%
}
\define@key{psset}{plane}{%			xy,xz,yz
	\edef\psk@ThreeDplot@plane{#1}%
}
\define@key{psset}{origin}{%			combination of (lr)(tBb)
	\edef\psk@ThreeDplot@origin{#1}%
}
\def\ThreeDplot@planeXY{xy}
\def\ThreeDplot@planeXZ{xz}
\def\ThreeDplot@planeYZ{yz}
%
% -------------- the length and node definitions -------------
%
\define@key{psset}{XO}{%			the X-offset
	\edef\psk@ThreeDplot@XO{#1}%
}
\define@key{psset}{YO}{%			the y-offset
	\edef\psk@ThreeDplot@YO{#1}%
}
\define@key{psset}{angleStep}{%		for circles
	\edef\psk@ThreeDplot@angleStep{#1}%
}
\define@key{psset}{posStart}{%		where the arrows start
	\edef\psk@ThreeDplot@posStart{#1}%
}
\define@key{psset}{length}{%		the length of the before|outlines
	\edef\psk@ThreeDplot@length{#1}%
}
\define@key{psset}{arrowOffset}{%	offset for \arrowLine
	\edef\psk@ThreeDplot@arrowOffset{#1}%
}
\define@key{psset}{visibleLineStyle}{%	offset for \arrowLine
	\edef\psk@ThreeDplot@visibleLineStyle{#1}%
}
\define@key{psset}{invisibleLineStyle}{%	offset for \arrowLine
	\edef\psk@ThreeDplot@invisibleLineStyle{#1}%
}
%
\define@key{psset}{nameX}{%		start of the object arrow
	\edef\psk@ThreeDplot@nameX{#1}%
}
\define@key{psset}{spotX}{%		where to draw the label
	\edef\psk@ThreeDplot@spotX{#1}%
}
\define@key{psset}{nameY}{%
	\edef\psk@ThreeDplot@nameY{#1}%
}
\define@key{psset}{spotY}{%
	\edef\psk@ThreeDplot@spotY{#1}%
}
\define@key{psset}{nameZ}{%
	\edef\psk@ThreeDplot@nameZ{#1}%
}
\define@key{psset}{spotZ}{%
	\edef\psk@ThreeDplot@spotZ{#1}%
}
%
\def\psds@none{\pst@gdot{}}% define none for the dotstyle
%
%%%%%%%%%%%%%%%%%%%%%%%%%%%%%%%%%%%%%%%%%%%%%%%%%%%%%%%%%%%%%%%%
%                                                              %
%                        B A S I C S	                       %
%                                                              %
%%%%%%%%%%%%%%%%%%%%%%%%%%%%%%%%%%%%%%%%%%%%%%%%%%%%%%%%%%%%%%%%
%
\pstheader{pst-3dplot.pro}
\def\pst@3ddict{tx@3DPlotDict begin }
\def\tx@saveCoor{\pst@3ddict saveCoor end }
\def\tx@ConvertTo2D{\pst@3ddict ConvertTo2D end }
\def\tx@SphericalCoor{\pst@3ddict SphericalCoor end }
\def\tx@ConvertToCartesian{\pst@3ddict ConvertToCartesian end }
%
\def\setDefaults{%
    \setkeys{psset}{%
	drawing=true,hiddenLine=false,%
	xMin=-1,xMax=4,yMin=-1,yMax=4,zMin=-1,zMax=4,%
	xThreeDunit=1,yThreeDunit=1,zThreeDunit=1,%
	Alpha=45,Beta=30,%
	drawStyle=xLines,%
	xPlotpoints=25,yPlotpoints=25,%
	beginAngle=0,endAngle=360,%
	linejoin=1,%
	XO=0,YO=0,%
	angleStep=1,%
	posStart=0,length=2,%
	arrowOffset=0,%
	visibleLineStyle=solid,%
	invisibleLineStyle=dashed,%
	nameX=x,spotX=180,%
	nameY=y,spotY=0,%
	nameZ=z,spotZ=90,%
	plane=xy,origin=c,
	SphericalCoor=false%
    }%
}
\setDefaults
%
% (#1) -> #1 #2 #3
\def\getThreeDCoor#1#2{\pst@getThreeDcoor(#1)\let#2\pst@coorIIID}
\def\pst@getThreeDcoor(#1,#2,#3){%
	\edef\pst@coorIIID{#1\space #2\space #3\space}
}
\def\pst@addThreeDVec(#1)(#2)#3#4#5{%
	\pst@calcThreeDVec(#1)(#2)%
	\let#3\pst@getValueX%
	\let#4\pst@getValueY%
	\let#5\pst@getValueZ%
}
\def\pst@subThreeDVec(#1)(#2,#3,#4)#5#6#7{%
	\pst@calcThreeDVec(#1)(-#2,-#3,-#4)%
	\let#5\pst@getValueX%
	\let#6\pst@getValueY%
	\let#7\pst@getValueZ%
}
\def\pst@calcThreeDVec(#1,#2,#3)(#4,#5,#6){%
  \begingroup
    \pst@dima=#1 pt%
    \pst@dimb=#2 pt%
    \pst@dimc=#3 pt%
	\pst@dimf=#4 pt%
	\pst@dimg=#5 pt%
	\pst@dimh=#6 pt%
	\advance\pst@dima by \pst@dimf%
	\advance\pst@dimb by \pst@dimg%
	\advance\pst@dimc by \pst@dimh%
	\xdef\pst@getValueX{\pst@number\pst@dima}%
	\xdef\pst@getValueY{\pst@number\pst@dimb}%
	\xdef\pst@getValueZ{\pst@number\pst@dimc}%
  \endgroup
}%
%
% the ThreeD coordinate system
%
\def\pstThreeDCoor{\@ifnextchar[{\pst@ThreeDCoor}{\pst@ThreeDCoor[]}}
\def\pst@ThreeDCoor[#1]{{%
	\pst@killglue
	\setkeys{psset}{linewidth=0.5pt,linecolor=red,arrows=->,dotstyle=|}
	\setkeys{psset}{#1}
	\pstThreeDNode(\psk@ThreeDplot@xMin,0,0){xMin}
	\pstThreeDNode(\psk@ThreeDplot@xMax,0,0){xMax}
	\pstThreeDNode(0,\psk@ThreeDplot@yMin,0){yMin}
	\pstThreeDNode(0,\psk@ThreeDplot@yMax,0){yMax}
	\pstThreeDNode(0,0,\psk@ThreeDplot@zMin){zMin}
	\pstThreeDNode(0,0,\psk@ThreeDplot@zMax){zMax}
	\ifPst@drawing
	    % ThreeDplot axes
	    \psline(xMin)(xMax)
	    \psline(yMin)(yMax)
	    \psline(zMin)(zMax)
	    \uput[\psk@ThreeDplot@spotX](xMax){$\mathbf{\psk@ThreeDplot@nameX}$}
	    \uput[\psk@ThreeDplot@spotY](yMax){$\mathbf{\psk@ThreeDplot@nameY}$}
	    \uput[\psk@ThreeDplot@spotZ](zMax){$\mathbf{\psk@ThreeDplot@nameZ}$}
	\fi% ifdrawing
}}
%
% put anything at (#2,#3,#4)
%
\def\pstThreeDPut{\@ifnextchar[{\pst@ThreeDPut}{\pst@ThreeDPut[]}}
\def\pst@ThreeDPut[#1](#2,#3,#4)#5{{%
	\pst@killglue%
	\setkeys{psset}{#1}
	\pstThreeDNode(#2,#3,#4){temp@pstNode}%
%	\def\@tempa{c}
%	\ifx\pst@ThreeDplot@origin\@tempa%
%	    \rput(temp@pstNode){#5}%
%	\else%
	    \rput[\psk@ThreeDplot@origin](temp@pstNode){#5}%
%	\fi%
}}
%
% draws a 3d line
%
\def\cartesianIIID@coor#1,#2,#3,#4\@nil{\edef\pst@coor{#1 #2 #3 }}
\def\NormalIIIDCoor{%
    \def\pst@@getcoor##1{\pst@expandafter\cartesianIIID@coor{##1},\relax,\@nil}%
    \def\pst@@getangle##1{%
        \pst@checknum{##1}\pst@angle
        \edef\pst@angle{\pst@angle \pst@angleunit}%
    }%
    \def\psput@##1{\pst@@getcoor{##1}\leavevmode\psput@cartesian}%
}
\def\pstThreeDLine{\NormalIIIDCoor\pst@object{lineIIID}}
\def\lineIIID@i{%
    \pst@getarrows{%
        \begin@OpenObj
        \pst@getcoors[\lineIIID@ii%
    }%
}
\def\lineIIID@ii{%
    \addto@pscode{%
		\pst@3ddict
		/Alpha \psk@ThreeDplot@Alpha\space def %
		/Beta \psk@ThreeDplot@Beta\space def %
		/SphericalCoor {\ifPst@SphericalCoor true \else false \fi} def %
		/xUnit { \pst@number\psxunit\space mul } def
		/yUnit { \pst@number\psyunit\space mul } def
		convertStackTo2D
		end
		\pst@cp \psline@iii \tx@Line%
	}%
    \end@OpenObj%
	\SpecialCoor%
}
%
% set a 3d dot
%
\def\pstThreeDDot{\@ifnextchar[{\pst@ThreeDDot}{\pst@ThreeDDot[]}}
\def\pst@ThreeDDot[#1](#2,#3,#4){{%
%	\pst@killglue
	\setkeys{psset}{linestyle=dashed,linewidth=0.5pt}% default
	\setkeys{psset}{#1}
	\pstThreeDNode[#1](#2,#3,#4){A}
	\ifx\psk@dotstyle\@none\else\psdots(A)\fi
	\ifPst@drawCoor
		\pnode(!%
			\pst@3ddict
			/Alpha \psk@ThreeDplot@Alpha\space def
			/Beta \psk@ThreeDplot@Beta\space def
			\ifPst@SphericalCoor
				#2\space #3\space #4\space
				ConvertToCartesian /z 0 def
			\else
				/x #2 def /y #3 def /z 0 def
			\fi
		ConvertTo2D x2D y2D end ){B}
		\psline(A)(B)
		\pnode(!%   y=0, z=0
			\pst@3ddict
			\ifPst@SphericalCoor
				#2\space #3\space #4\space
				ConvertToCartesian /y 0 def /z 0 def
			\else
				/x #2 def /y 0 def /z 0 def
			\fi
		ConvertTo2D x2D y2D end ){A}%
		\psline(B)(A)%
		\pnode(!%   x=0, z=0
			\pst@3ddict
			\ifPst@SphericalCoor
				#2\space #3\space #4\space
				ConvertToCartesian /x 0 def /z 0 def
			\else
				/x 0 def /y #3 def /z 0 def
			\fi
		ConvertTo2D x2D y2D end ){A}%
		\psline(B)(A)
	\fi
}}
%
% transform the 3d coordinates of the node (#1,#2,#3)
% into a 2d node with the name #4
%
\def\pstThreeDNode{\@ifnextchar[{\pst@ThreeDNode}{\pst@ThreeDNode[]}}
\def\pst@ThreeDNode[#1](#2,#3,#4)#5{{%
	\psset{#1}
	\pnode(!%
		\pst@3ddict
		/Alpha \psk@ThreeDplot@Alpha\space def
		/Beta \psk@ThreeDplot@Beta\space def
		\ifPst@SphericalCoor
			#2\space #3\space #4\space
			ConvertToCartesian
		\else
			/x #2 def /y #3 def /z #4 def
		\fi
		ConvertTo2D x2D y2D end){#5}%
}}
%
% a 3d uput[](){}
%
\def\pstUThreeDPut{\@ifnextchar[{\pst@UThreeDPut}{\pst@UThreeDPut[]}}
\def\pst@UThreeDPut[#1](#2,#3,#4)#5{{%
	\uput[#1](!%
		\pst@3ddict
		/Alpha \psk@ThreeDplot@Alpha\space def
		/Beta \psk@ThreeDplot@Beta\space def
		\ifPst@SphericalCoor
			#2\space #3\space #4\space
			ConvertToCartesian
		\else
			/x #2 def /y #3 def /z #4 def
		\fi
		ConvertTo2D x2D y2D end){#5}
}}
%
% Trangle [options](Point1)(Point2)(Point3)
%
\def\pstThreeDTriangle{%
	\@ifnextchar[%
		{\do@ThreeDTriangle}%
		{\do@ThreeDTriangle[]}%
}
\def\do@ThreeDTriangle[#1](#2)(#3)(#4){{%
	\ifx#1\empty\else\setkeys{psset}{#1}\fi%
	\ifPst@drawCoor
		\pstThreeDDot[drawCoor=true](#2)%
		\pstThreeDDot[drawCoor=true](#3)%
		\pstThreeDDot[drawCoor=true](#4)%
	\fi
	\pstThreeDNode(#2){A}
	\pstThreeDNode(#3){B}
	\pstThreeDNode(#4){C}
	\ifx\psk@fillstyle\@none
	\else
		\pscustom{%
			\code{\psk@ThreeDplot@linejoin\space setlinejoin}%
	\fi%
		\psline[#1](A)(B)(C)(A)
	\ifx\psk@fillstyle\@none\else}\fi%
}}
%
% draws a threeD square as a polygon
%
% [#1] options
% (#2) starting vector ax,ay,az
% (#3) first direction vector ux,uy,uz
% (#4) second direction vector wx,wy,wz
%
\def\pstThreeDSquare{\pst@object{squareIIID}}
\def\squareIIID@i{\squareIIID@ii}
\def\squareIIID@ii(#1)(#2)(#3){%
	\ifPst@drawCoor%
		{%
		\psset{linestyle=dashed,linewidth=0.5pt}%
		\pstThreeDDot(#1)%
		\pst@addThreeDVec(#1)(#2)\pst@tempa\pst@tempb\pst@tempc%
		\pstThreeDDot(\pst@tempa,\pst@tempb,\pst@tempc)%
		\pst@addThreeDVec(#1)(#3)\pst@tempa\pst@tempb\pst@tempc
		\pstThreeDDot(\pst@tempa,\pst@tempb,\pst@tempc)%
		\pst@addThreeDVec(\pst@tempa,\pst@tempb,\pst@tempc)%
			(#2)\pst@tempa\pst@tempb\pst@tempc%
		\pstThreeDDot(\pst@tempa,\pst@tempb,\pst@tempc)%
		}%
	\fi%
	\getThreeDCoor{#1}\pst@tempa%
	\getThreeDCoor{#2}\pst@tempb%
	\getThreeDCoor{#3}\pst@tempc%
	\begin@OpenObj
		\addto@pscode{%
			/P1 {\pst@tempa } def % x y z or Radius longitude lattitude
			/P2 {\pst@tempb } def %
			/P3 {\pst@tempc } def %
			/Alpha \psk@ThreeDplot@Alpha\space def %
			/Beta \psk@ThreeDplot@Beta\space def %
			/SphericalCoor \ifPst@SphericalCoor true \else false \fi def %
			/xUnit { \pst@number\psxunit\space mul } def
			/yUnit { \pst@number\psyunit\space mul } def
			\pst@3ddict
			P1 saveCoor
			SphericalCoor {ConvertToCartesian} if
			ConvertTo2D
			/x0 x2D xUnit def /y0 y2D yUnit def
			P2 saveCoor
			SphericalCoor {ConvertToCartesian} if
			ConvertTo2D
			/x1 x2D xUnit x0 add def /y1 y2D yUnit y0 add def
			P3 saveCoor
			SphericalCoor {ConvertToCartesian} if
			ConvertTo2D
			/x2 x2D xUnit x1 add def /y2 y2D yUnit y1 add def
			P2 saveCoor
			SphericalCoor {ConvertToCartesian} if
			ConvertTo2D
			/x3 x2D xUnit neg x2 add def /y3 y2D yUnit neg y2 add def
			[ x0 y0 x1 y1 x2 y2 x3 y3 x0 y0
			\pst@cp \psline@iii \tx@Line
			end
		}%
    \end@OpenObj%
}

%
% draws a threeD Box
% [#1] options
% (#2) first direction vector ux,uy,uz
% (#3) second direction vector vx,vy,vz
% (#4) third direction vector wx,wy,wz
%
\def\pstThreeDBox{\@ifnextchar[{\doThreeDBox}{\doThreeDBox[]}}
\def\doThreeDBox[#1](#2)(#3)(#4)(#5){{
	\psset{linestyle=\psk@ThreeDplot@invisibleLineStyle}
	\psset{#1}
	%
%	\psset{linestyle=dashed}
%	\psset{linestyle=\psk@ThreeDplot@invisibleLineStyle}
	\addto@par{linestyle=\psk@ThreeDplot@invisibleLineStyle}
	\pstThreeDSquare[#1](#2)(#4)(#5) % 	lower square
	\pstThreeDSquare(#2)(#3)(#4)% 	back square
	%
	\psset{linestyle=\psk@ThreeDplot@visibleLineStyle}
	\pst@addThreeDVec(#2)(#4)\pst@tempa\pst@tempb\pst@tempc
	\pstThreeDSquare(\pst@tempa,\pst@tempb,\pst@tempc)(#3)(#5)% left square
	\pst@addThreeDVec(#2)(#3)\pst@tempa\pst@tempb\pst@tempc
	\pstThreeDSquare(\pst@tempa,\pst@tempb,\pst@tempc)(#4)(#5)% top square
	\pst@addThreeDVec(#2)(#5)\pst@tempa\pst@tempb\pst@tempc
	\pstThreeDSquare(\pst@tempa,\pst@tempb,\pst@tempc)(#3)(#4)% front square
	%
}}
%
% Sphere
%
\def\pstThreeDSphere{\@ifnextchar[{\pst@ThreeDSphere}{\pst@ThreeDSphere[]}}
\def\pst@ThreeDSphere[#1](#2,#3,#4)#5{{%
	\psset{THETA=\psk@ThreeDplot@Beta,PHI=\psk@ThreeDplot@Beta,%
		Dobs=10,Decran=10}
	\psset{#1}
	\pstThreeDNode(#2,#3,#4){SphereCenter}
	\rput(SphereCenter){\SphereThreeD(0,0,0){#5}}
}}
%
% set a 3d ellipse/circle
%
% #1 options
% #2 center cx,cy,cz
% #3 radius ax,ay,az
% #4 radius bx,by,bz
%
\def\pstThreeDEllipse{\@ifnextchar[{\doThreeDEllipse}{\doThreeDEllipse[]}}
\def\doThreeDEllipse[#1](#2)(#3)(#4){{%
	\psset{plotstyle=curve,xPlotpoints=\psk@plotpoints}
	\psset{#1}
	\getThreeDCoor{#2}\pst@tempc%	 center
	\getThreeDCoor{#3}\pst@tempa%	 a
	\getThreeDCoor{#4}\pst@tempb%	 b
	\parametricplotThreeD(\psk@ThreeDplot@beginAngle,\psk@ThreeDplot@endAngle){%
		\pst@3ddict
		\pst@tempc\space
		\ifPst@SphericalCoor
			ConvertToCartesian x y z
		\fi
		\pst@tempa\space
		\ifPst@SphericalCoor
			ConvertToCartesian x y z
		\fi
		\pst@tempb\space
		\ifPst@SphericalCoor
			ConvertToCartesian x y z
		\fi
		9 -1 roll t cos 8 -1 roll neg mul add t sin 5 -1 roll mul add
		7 -1 roll t cos 7 -1 roll neg mul add t sin 5 -1 roll mul add
		5 -1 roll t cos 6 -1 roll neg mul add t sin 5 -1 roll mul add
		end
	}
}}
%
\def\pstThreeDCircle{\pstThreeDEllipse}
\def\pstThreeDPlotFunc{\psplotThreeD}%	only for compatibility
%
\def\pstRotPoint{\@ifnextchar[{\pstRotPointi}{\pstRotPointi[]}}
\def\pstRotPointi[#1](#2,#3,#4)(#5,#6,#7)#8{{%
	\setkeys{psset}{#1}
	\def\pst@ThreeDplot@ThetaX{#5}	% rotating angle
	\def\pst@ThreeDplot@ThetaY{#6}	% rotating angle
	\def\pst@ThreeDplot@ThetaZ{#7}	% rotating angle
	\pnode(!%
		\pst@3ddict
		/Alpha \psk@ThreeDplot@Alpha\space def
		/Beta \psk@ThreeDplot@Beta\space def
		\ifPst@SphericalCoor
			#2\space #3\space #4\space
			ConvertToCartesian
		\else
			/x #2 def /y #3 def /z #4 def
		\fi
% rotating on the z-Axes
		/Theta \pst@ThreeDplot@ThetaZ\space def
		/x0 x Theta cos mul y Theta sin mul sub def
		/y x Theta sin mul y Theta cos mul add def
		/x x0 def
% rotating on the y-Axes
		/Theta \pst@ThreeDplot@ThetaY\space def
		/x0 x Theta cos mul z Theta sin mul add def
		/z x neg Theta sin mul z Theta cos mul add def
		/x x0 def
% rotating on the x-Axes
		/Theta \pst@ThreeDplot@ThetaX\space def
		/y0 y Theta cos mul z Theta sin mul sub def
		/z y Theta sin mul z Theta cos mul add def
		/y y0 def
%
		ConvertTo2D x2D y2D end){#8}
}}
%
%%%%%%%%%%%%%%%%%%%%%%%%%%%%%%%%%%%%%%%%%%%%%%%%%%%%%%%%%%%%%%%%%%%%%%%%%%%%
%
\def\psplotinit#1{\xdef\psplot@init{#1 }}
\def\psplot@init{}
%
\def\psplotThreeD{\@ifnextchar[{\ps@plotThreeD}{\ps@plotThreeD[]}}
\def\ps@plotThreeD[#1](#2,#3)(#4,#5)#6{%
    \setkeys{psset}{hiddenLine=false,drawStyle=xLines}% default
    \setkeys{psset}{#1}
    \ifx\psk@ThreeDplot@drawStyle\drawStyle@xyLines
	\doPlotThreeD@xLines(#2,#3)(#4,#5){#6}
	\doPlotThreeD@yLines(#2,#3)(#4,#5){#6}
    \else
	\ifx\psk@ThreeDplot@drawStyle\drawStyle@yxLines
	    \doPlotThreeD@yLines(#2,#3)(#4,#5){#6}
	    \doPlotThreeD@xLines(#2,#3)(#4,#5){#6}
	\else
	    \ifx\psk@ThreeDplot@drawStyle\drawStyle@xLines
		\doPlotThreeD@xLines(#2,#3)(#4,#5){#6}
	    \fi
	    \ifx\psk@ThreeDplot@drawStyle\drawStyle@yLines
		\doPlotThreeD@yLines(#2,#3)(#4,#5){#6}
	    \fi
	\fi
    \fi
}
\def\doPlotThreeD@xLines{%
	\def\pst@par{}
	\pst@object{psplotThreeD@xLines}%
}
\def\doPlotThreeD@yLines{%
	\def\pst@par{}
	\pst@object{psplotThreeD@yLines}%
}
\def\psplotThreeD@xLines@i(#1,#2)(#3,#4)#5{%
    \pst@killglue
    \begingroup
	\use@par%
	\@pstfalse
	\@nameuse{beginplot@\psplotstyle}%
	\gdef\psplot@init{}%
	\@nameuse{testqp@\psplotstyle}%
	\if@pst%	lines and dots
	    \addto@pscode{%
		/Alpha \psk@ThreeDplot@Alpha\space def
		/Beta \psk@ThreeDplot@Beta\space def
		/xMin #1 def
		/x xMin def
		/x1 #2 def
		/y #3 def
		/y1 #4 def
		/dx x1 x sub \psk@ThreeDplot@xPlotpoints\space div def
		/dy y1 y sub \psk@ThreeDplot@yPlotpoints\space div def
		/func {#5} def
		/xyz {
		    x neg Alpha cos mul y Alpha sin mul add \pst@number\psxunit mul
		    x Alpha sin mul y Alpha cos mul add neg Beta sin mul
			func Beta cos mul add \pst@number\psyunit mul} def
	    }%
	    \psplotThreeD@xLines@ii
	\else%	curves
	    \endgroup
	    \multido{\n@Y=0+1}{\psk@ThreeDplot@yPlotpoints}{%
		\ifPst@hiddenLine\pscustom[fillcolor=white,fillstyle=solid]{\fi
		    \@nameuse{beginplot@\psplotstyle}%
			\addto@pscode{%
			    /Alpha \psk@ThreeDplot@Alpha\space def
			    /Beta \psk@ThreeDplot@Beta\space def
			    /xMin #1 def
			    /x xMin def
			    /x1 #2 def
			    /dx x1 x sub \psk@ThreeDplot@xPlotpoints\space div def
			    /dy #4\space #3\space sub \psk@ThreeDplot@yPlotpoints\space div def
			    /y #3\space \n@Y\space dy mul add def
			    /func {#5} def
			    /xyz {
				x neg Alpha cos mul y Alpha sin mul add \pst@number\psxunit mul
				x Alpha sin mul y Alpha cos mul add neg Beta sin mul
				    func Beta cos mul add \pst@number\psyunit mul} def
			}%
		    \psplotThreeD@xLines@iii%
		\ifPst@hiddenLine }\fi%
	    }%
	\fi
    \endgroup
    \ignorespaces%
}
\def\psplotThreeD@xLines@ii{%
    \addto@pscode{%
	xyz \@nameuse{beginqp@\psplotstyle}
	\psk@ThreeDplot@yPlotpoints\space 1 add {
	    /x xMin def
	    xyz moveto
	    \psk@ThreeDplot@xPlotpoints\space 1 add {
		xyz \@nameuse{doqp@\psplotstyle}
		/x x dx add def
	    } repeat
	    /y y dy add def
	} repeat
	/x x1 def
	/y y1 def
	xyz \@nameuse{doqp@\psplotstyle}
    }%
    \@nameuse{endqp@\psplotstyle}%
}
\def\psplotThreeD@xLines@iii{%    curves
    \addto@pscode{%
	    /x xMin def
	    mark
	    /n 2 def
		\psk@ThreeDplot@xPlotpoints {
	    	    xyz
    		    n 2 roll
    		    /n n 2 add def
    		    /x x dx add def
    		} repeat
        /x x1 def
        xyz
        n 2 roll
    }%
    \@nameuse{endplot@\psplotstyle}%
}
%
\def\psplotThreeD@yLines@i(#1,#2)(#3,#4)#5{%
    \pst@killglue
    \begingroup
	\use@par%
	\@pstfalse
	\@nameuse{beginplot@\psplotstyle}%
	\gdef\psplot@init{}%
	\@nameuse{testqp@\psplotstyle}%
	\if@pst%	lines and dots
	    \addto@pscode{%
		/Alpha \psk@ThreeDplot@Alpha\space def
		/Beta \psk@ThreeDplot@Beta\space def
    		/x #1 def
		/x1 #2 def
		/yMin #3 def
		/y yMin def
		/y1 #4 def
		/dx x1 x sub \psk@ThreeDplot@xPlotpoints\space div def
		/dy y1 y sub \psk@ThreeDplot@yPlotpoints\space div def
		/func {#5} def
		/xyz {
		    x neg Alpha cos mul y Alpha sin mul add \pst@number\psxunit mul
		    x Alpha sin mul y Alpha cos mul add neg Beta sin mul
			func Beta cos mul add \pst@number\psyunit mul} def
	    }%
	    \psplotThreeD@yLines@ii
	\else%	curves
	    \endgroup
	    \multido{\n@X=0+1}{\psk@ThreeDplot@xPlotpoints}{%
		\ifPst@hiddenLine\pscustom[fillcolor=white,fillstyle=solid]{\fi
		    \@nameuse{beginplot@\psplotstyle}%
			\addto@pscode{%
			    /Alpha \psk@ThreeDplot@Alpha\space def
			    /Beta \psk@ThreeDplot@Beta\space def
			    /yMin #3 def
			    /y yMin def
			    /y1 #4 def
			    /dy y1 y sub \psk@ThreeDplot@yPlotpoints\space div def
			    /dx #2\space #1\space sub \psk@ThreeDplot@xPlotpoints\space div def
			    /x #1\space \n@X\space dx mul add def
			    /func {#5} def
			    /xyz {
				x neg Alpha cos mul y Alpha sin mul add \pst@number\psxunit mul
				x Alpha sin mul y Alpha cos mul add neg Beta sin mul
				    func Beta cos mul add \pst@number\psyunit mul} def
			}%
		    \psplotThreeD@yLines@iii%
		\ifPst@hiddenLine }\fi%
	    }%
	\fi
    \endgroup
    \ignorespaces%
}
\def\psplotThreeD@yLines@ii{%
    \addto@pscode{%
	xyz \@nameuse{beginqp@\psplotstyle}
	\psk@ThreeDplot@xPlotpoints\space 1 add {
	    /y yMin def
	    xyz moveto
	    \psk@ThreeDplot@yPlotpoints\space 1 add {
		xyz \@nameuse{doqp@\psplotstyle}
		/y y dy add def
	    } repeat
	    /x x dx add def
	} repeat
	/x x1 def
	/y y1 def
	xyz \@nameuse{doqp@\psplotstyle}
    }%
    \@nameuse{endqp@\psplotstyle}%
}
\def\psplotThreeD@yLines@iii{%    curves
    \addto@pscode{%
	    /y yMin def
	    mark
	    /n 2 def
		\psk@ThreeDplot@yPlotpoints {
	    	    xyz
    		    n 2 roll
    		    /n n 2 add def
    		    /y y dy add def
    		} repeat
        /y y1 def
        xyz
        n 2 roll
    }%
    \@nameuse{endplot@\psplotstyle}%
}
%
% parametricplot ----------------------------------------------------------------
%
\def\parametricplotThreeD{%
	\@ifnextchar[{\doParametricPlotThreeD}{\doParametricPlotThreeD[]}%
}
\def\doParametricPlotThreeD[#1]{%
	\setkeys{psset}{#1}%
	\def\pst@par{}%
	\pst@object{psParamPlotThreeD}%
}
\def\psParamPlotThreeD@i(#1){%
	\@ifnextchar({\psParamPlotThreeD@ii(#1)}{\psParamPlotThreeD@ii(#1)(0,0)}%
}
\def\psParamPlotThreeD@ii(#1,#2)(#3,#4)#5 {%
    \pst@killglue
    \begingroup
	\use@par%
	\@pstfalse
	\@nameuse{beginplot@\psplotstyle}%
	\gdef\parametricII@init{}%
	\@nameuse{testqp@\psplotstyle}%
	\if@pst
	    \addto@pscode{%
	    	/Alpha \psk@ThreeDplot@Alpha\space def
	    	/Beta \psk@ThreeDplot@Beta\space def
		/tMin #1 def
		/t tMin def
		/t1 #2 def
		/u #3 def
		/u1 #4 def
		/dt t1 t sub \psk@ThreeDplot@xPlotpoints\space div def
		/du u1 u sub \psk@ThreeDplot@yPlotpoints\space div def
		/xyz {
		    #5
		    3 -1 roll               % y z x
		    \psk@ThreeDplot@xThreeDunit\space div
		    dup                     % y z x x
		    neg Alpha cos mul       % y z x (x)
		    4 -1 roll               % z x (x) y
		    \psk@ThreeDplot@yThreeDunit\space div
		    dup                     % z x (x) y y
		    5 1 roll                % y z x (x) y
		    Alpha sin mul add \pst@number\psxunit mul       % y z x (x2)
		    4 1 roll                % (x2) y z x
		    Alpha sin mul           % (x2) y z (x)
		    3 -1 roll               % (x2) z (x) y
		    Alpha cos mul add neg Beta sin mul % (x2) z (x)
		    exch                    % (x2) (x) z
		    \psk@ThreeDplot@zThreeDunit\space div
		    Beta cos mul add \pst@number\psyunit mul        % (x2) (y2)
		} def
	    }%
	    \psParamPlotThreeD@iii%
	\else
	    \endgroup
	    \ifx#3#4\psset{yPlotpoints=1}\fi% do we have two different parameters or #3=#4
		\multido{\n@Y=0+1}{\psk@ThreeDplot@yPlotpoints}{%
		\@nameuse{beginplot@\psplotstyle}%
		\addto@pscode{%
		    /Alpha \psk@ThreeDplot@Alpha\space def
		    /Beta \psk@ThreeDplot@Beta\space def
		    /tMin #1 def
		    /t tMin def
		    /t1 #2 def
		    /dt t1 t sub \psk@ThreeDplot@xPlotpoints\space div def
		    /du #4\space #3\space sub \psk@ThreeDplot@yPlotpoints\space div def
		    /u #3\space \n@Y\space du mul add def
		    /xyz {
			#5
			3 -1 roll               % y z x
			\psk@ThreeDplot@xThreeDunit\space div
			dup                     % y z x x
			neg Alpha cos mul       % y z x (x)
			4 -1 roll               % z x (x) y
			\psk@ThreeDplot@yThreeDunit\space div
			dup                     % z x (x) y y
			5 1 roll                % y z x (x) y
			Alpha sin mul add \pst@number\psxunit mul       % y z x (x2)
			4 1 roll                % (x2) y z x
			Alpha sin mul           % (x2) y z (x)
			3 -1 roll               % (x2) z (x) y
			Alpha cos mul add neg Beta sin mul % (x2) z (x)
			exch                    % (x2) (x) z
			\psk@ThreeDplot@zThreeDunit\space div
			Beta cos mul add \pst@number\psyunit mul        % (x2) (y2)
		    } def
		}%
		\psParamPlotThreeD@iv%
	    }%
	\fi
    \endgroup
    \ignorespaces%
}
\def\psParamPlotThreeD@iii{%
    \addto@pscode{%
	xyz \@nameuse{beginqp@\psplotstyle}
	\psk@ThreeDplot@yPlotpoints\space 1 add {
	    /t tMin def
	    xyz moveto
	    \psk@ThreeDplot@xPlotpoints\space 1 add {
		xyz \@nameuse{doqp@\psplotstyle}
		/t t dt add def
	    } repeat
	    /u u du add def
	} repeat
	/t u1 def
	/u u1 def
	xyz \@nameuse{doqp@\psplotstyle}
    }%
    \@nameuse{endqp@\psplotstyle}%
}
\def\psParamPlotThreeD@iv{%
	\addto@pscode{%
%	    /t tMin def
		mark
	    /n 2 def
		\psk@ThreeDplot@xPlotpoints {
			xyz
    		n 2 roll
    		/n n 2 add def
    		/t t dt add def
		} repeat
		/t t1 def
		xyz
		n 2 roll
    }%
	\@nameuse{endplot@\psplotstyle}%
}
%
% Plot 3D data
%
\def\fileplotThreeD{\def\pst@par{}\pst@object{fileplotThreeD}}
\def\fileplotThreeD@i#1{%
    \pst@killglue
    \begingroup
	\use@par
	\@pstfalse
	\@nameuse{testqp@\psplotstyle}%
	\if@pst
	    \dataplotThreeD@ii{\pst@readfile{#1}}%
	\else
	    \listplotThreeD@ii{\pst@altreadfile{#1}}%
	\fi
    \endgroup
    \ignorespaces%
}
%
% dataplot
%
\def\dataplotThreeD{\def\pst@par{}\pst@object{dataplotThreeD}}
\def\dataplotThreeD@i#1{%
  \pst@killglue
  \begingroup
    \use@par
    \@pstfalse
    \@nameuse{testqp@\psplotstyle}%
    \if@pst
      \dataplotThreeD@ii{\addto@pscode{#1}}%
    \else
      \listplot@ii{\addto@pscode{#1}}%
    \fi
  \endgroup
  \ignorespaces%
}
\def\dataplotThreeD@ii#1{%
  \@nameuse{beginplot@\psplotstyle}%
    \addto@pscode{%
      /Alpha \psk@ThreeDplot@Alpha\space def
      /Beta \psk@ThreeDplot@Beta\space def
      /Dx { /D { Dy } def } def
      /Dy { /D { Dz } def } def
      /Dz { 			% now we have x y z
	3 -1 roll		% y z x
        \psk@ThreeDplot@xThreeDunit\space div
	dup			% y z x x
	neg Alpha cos mul	% y z x (x)
	4 -1 roll		% z x (x) y
	\psk@ThreeDplot@yThreeDunit\space div
	dup			% z x (x) y y
	5 1 roll		% y z x (x) y
	Alpha sin mul add \pst@number\psxunit mul % y z x (x2)
	4 1 roll		% (x2) y z x
	Alpha sin mul		% (x2) y z (x)
	3 -1 roll		% (x2) z (x) y
	Alpha cos mul add neg Beta sin mul % (x2) z (x)
	exch			% (x2) (x) z
	\psk@ThreeDplot@zThreeDunit\space div
	Beta cos mul add \pst@number\psyunit mul % (x2) (y2)
	Do /D { Dx } def } def
      /D { /D { Dx } def } def
      /Do {
        \@nameuse{beginqp@\psplotstyle}%
        /Do { \@nameuse{doqp@\psplotstyle}} def
      } def}%
    #1%
    \addto@pscode{D}%
  \@nameuse{endqp@\psplotstyle}%
}
%
%
\pst@def{ScalePointsThreeD}<%
	counttomark dup dup cvi eq not { exch pop } if
	/m exch def /n m 3 div cvi def
	n { 				% now we have x y z
		3 -1 roll		% y z x
		\psk@ThreeDplot@xThreeDunit\space div
		dup			% y z x x
		neg Alpha cos mul	% y z x (x)
		4 -1 roll		% z x (x) y
		\psk@ThreeDplot@yThreeDunit\space div
		dup			% z x (x) y y
		5 1 roll		% y z x (x) y
		Alpha sin mul add \pst@number\psxunit mul	% y z x (x2)
		4 1 roll		% (x2) y z x
		Alpha sin mul		% (x2) y z (x)
		3 -1 roll		% (x2) z (x) y
		Alpha cos mul add neg Beta sin mul % (x2) z (x)
		exch			% (x2) (x) z
		\psk@ThreeDplot@zThreeDunit\space div
		Beta cos mul add \pst@number\psyunit mul 	% (x2) (y2)
		m 1 sub 1 roll m 1 sub 1 roll /m m 3 sub def } repeat>
%
% listplotThreeD
%
\def\listplotThreeD{\def\pst@par{}\pst@object{listplotThreeD}}
\def\listplotThreeD@i#1{\listplotThreeD@ii{\addto@pscode{#1}}}
\def\listplotThreeD@ii#1{%
	\@nameuse{beginplot@\psplotstyle}%
	\addto@pscode{%
		/Alpha \psk@ThreeDplot@Alpha\space def
		/Beta \psk@ThreeDplot@Beta\space def
		/D {} def mark %
	}%
	#1%
	\addto@pscode{\tx@ScalePointsThreeD}%
	\@nameuse{endplot@\psplotstyle}%
}
%
% adopted from pst-3d
%
\pst@def{TMSave}<%
  tx@Dict /TMatrix known not { /TMatrix { } def /RAngle { 0 } def } if
  /TMatrix [ TMatrix CM ] cvx def>
%
\pst@def{TMRestore}<%
  CP /TMatrix [ TMatrix setmatrix ] cvx def moveto>
%
\pst@def{TMChange}<%
  \tx@TMSave
  /cp [ currentpoint ] cvx def
  CM
  CP T \tx@STV
  CM matrix invertmatrix 
  matrix concatmatrix
  exch exec
  concat cp moveto>
%
\def\pstPlanePut{\@ifnextchar[{\pst@PlanePut}{\pst@PlanePut[]}}
\def\pst@PlanePut[#1](#2)#3{{%
	\pst@killglue%
	\setkeys{psset}{origin=lB}%
	\setkeys{psset}{#1}%
	\pstThreeDPut[#1](#2){\ps@Plane{#3}}%
}}
%
\def\ps@Plane{\pst@makebox{\psPlane@}}
\def\psPlane@{%
    \begingroup
	\leavevmode%
    	\hbox{%
%	    \kern\wd\pst@hbox%
	    \pst@Verb{%
		{ [ 1 0
		    /Alpha \psk@ThreeDplot@Alpha\space def
		    /Beta \psk@ThreeDplot@Beta\space def
		    /Delta Beta sin Alpha sin mul Alpha cos atan neg 90 add def 
		    /Gamma Beta sin Alpha cos mul neg Alpha sin atan def 
		    \ifx\psk@ThreeDplot@plane\ThreeDplot@planeXY
			270 Delta sub rotate
			/Rho 90 Gamma add Delta add def
		    \else%
			\ifx\psk@ThreeDplot@plane\ThreeDplot@planeXZ
		    	    270 Delta sub rotate
 			    /Rho 180 Delta add def
			\else%
		    	    Gamma rotate
			    /Rho 90 Gamma sub def
			\fi%
		    \fi%
		    Rho cos Rho sin 0 0 ] concat} \tx@TMChange%
	    }%
	    \box\pst@hbox%
	    \pst@Verb{\tx@TMRestore}%
%	    \kern\ht\pst@hbox%
	}%
    \endgroup%
}
%
%%%%%%%%%%%%%%%%%%%%%%%%%%%%%%%%%%%%%%%%%%%%%%%%%%%%%%%%%%%%%%%%%
% Utility stuff
%%%%%%%%%%%%%%%%%%%%%%%%%%%%%%%%%%%%%%%%%%%%%%%%%%%%%%%%%%%%%%%%%
% posStart=Starting point
% length=  Arrow length.
\def\Arrows{\@ifnextchar[{\pst@Arrows}{\pst@Arrows[]}}
\def\pst@Arrows[#1](#2)(#3){{%
	\setkeys{psset}{#1}%
	\pst@getcoor{#2}\pst@tempa
	\pst@getcoor{#3}\pst@tempb
	\pnode(!%
		/StartArrow \psk@ThreeDplot@posStart\space def
		/LengthArrow \psk@ThreeDplot@length\space def
		\pst@tempa /YA exch \pst@number\psyunit div def
		/XA exch \pst@number\psxunit div def
		\pst@tempb /YB exch \pst@number\psyunit div def
		/XB exch \pst@number\psxunit div def
		/denominateur XB XA sub def
		/numerateur YB YA sub def
		/angleDirectionAB numerateur denominateur Atan def
		/XD StartArrow angleDirectionAB cos mul XA add def
		/YD StartArrow angleDirectionAB sin mul YA add def
		/XF XD LengthArrow angleDirectionAB cos mul add def
		/YF YD LengthArrow angleDirectionAB sin mul add def
		XD YD ){ArrowStart}
	\pnode(! XF YF){ArrowEnd}
	\psset{arrows=->}
	\psline[#1](ArrowStart)(ArrowEnd)%
}}
%
% draw a line (===) outside: #2-----#3=======#4
%
\def\psOutLine{\@ifnextchar[{\pst@ToDrawOut}{\pst@ToDrawOut[]}}
\def\pst@ToDrawOut[#1](#2)(#3)#4{{%
	\setkeys{psset}{#1}%
	\pst@getcoor{#2}\pst@tempa
	\pst@getcoor{#3}\pst@tempb
	\pnode(!%
		/LengthArrow \psk@ThreeDplot@length\space def
		\pst@tempa /YA exch \pst@number\psyunit div def
		/XA exch \pst@number\psxunit div def
		\pst@tempb /YB exch \pst@number\psyunit div def
		/XB exch \pst@number\psxunit div def
		/denominateur XB XA sub def
		/numerateur YB YA sub def
		/angleDirectionAB numerateur denominateur Atan def
		/Xend XB LengthArrow angleDirectionAB cos mul add def
		/Yend YB LengthArrow angleDirectionAB sin mul add def
		Xend Yend){#4}
	\psline[#1](#3)(#4)
}}
%
% draw a line (===) before: #4========#2-----#3
%
\def\psBeforeLine{\@ifnextchar[{\pst@BeforeLine}{\pst@BeforeLine[]}}
\def\pst@BeforeLine[#1](#2)(#3)#4{{%
	\setkeys{psset}{#1}%
	\pst@getcoor{#2}\pst@tempa
	\pst@getcoor{#3}\pst@tempb
	\pnode(!%
		/LengthArrow \psk@ThreeDplot@length\space def
		\pst@tempa /YA exch \pst@number\psyunit div def
		/XA exch \pst@number\psxunit div def
		\pst@tempb /YB exch \pst@number\psyunit div def
		/XB exch \pst@number\psxunit div def
		/denominateur XB XA sub def
		/numerateur YB YA sub def
		/angleDirectionAB numerateur denominateur Atan def
		/Xstart XA LengthArrow angleDirectionAB cos mul sub def
		/Ystart YA LengthArrow angleDirectionAB sin mul sub def
		Xstart Ystart){#4}
	\psline[#1](#4)(#2)%
}}
%
% intersection de deux droites
% 2 juillet 2001/ rewritten 2003-01-27 Herbert
%
\def\ABinterCD(#1)(#2)(#3)(#4)#5{%
    \pst@getcoor{#1}\pst@tempa
    \pst@getcoor{#2}\pst@tempb
    \pst@getcoor{#3}\pst@tempc
    \pst@getcoor{#4}\pst@tempd
\pnode(!%
    /YA \pst@tempa exch pop \pst@number\psyunit div def
    /XA \pst@tempa pop \pst@number\psxunit div def
    /YB \pst@tempb exch pop \pst@number\psyunit div def
    /XB \pst@tempb pop \pst@number\psxunit div def
    /YC \pst@tempc exch pop \pst@number\psyunit div def
    /XC \pst@tempc pop \pst@number\psxunit div def
    /YD \pst@tempd exch pop \pst@number\psyunit div def
    /XD \pst@tempd pop \pst@number\psxunit div def
    /dY1 YB YA sub def
    /dX1 XB XA sub def
    /dY2 YD YC sub def
    /dX2 XD XC sub def
    dX1 abs 0.01 lt {
	/m2 dY2 dX2 div def
	XA dup XC sub m2 mul YC add 
    }{
	dX2 abs 0.01 lt {
	    /m1 dY1 dX1 div def
	    XC dup XA sub m1 mul YA add 
	}{%
	    /m1 dY1 dX1 div def
	    /m2 dY2 dX2 div def
	    m1 XA mul m2 XC mul sub YA sub YC add m1 m2 sub div dup
	    XA sub m1 mul YA add 
	} ifelse
    } ifelse ){#5}
}
%
% draw a parallel line
%     #2---------#3
%         #4----------#5(new)  
\def\Parallel{\@ifnextchar[{\pst@Parallel}{\pst@Parallel[]}}
\def\pst@Parallel[#1](#2)(#3)(#4)#5{{%
	\setkeys{psset}{#1}%
	\pst@getcoor{#2}\pst@tempa
	\pst@getcoor{#3}\pst@tempb
	\pst@getcoor{#4}\pst@tempc
	\pnode(!%
		/LengthArrow \psk@ThreeDplot@length\space def
		\pst@tempa /YA exch \pst@number\psyunit div def
		/XA exch \pst@number\psxunit div def
		\pst@tempb /YB exch \pst@number\psyunit div def
		/XB exch \pst@number\psxunit div def
		\pst@tempc /YC exch \pst@number\psyunit div def
		/XC exch \pst@number\psxunit div def
		/denominateur XB XA sub def
		/numerateur YB YA sub def
		/angleDirectionAB numerateur denominateur Atan def
		/XstartParallel XC LengthArrow angleDirectionAB cos mul add def
		/YstartParallel YC LengthArrow angleDirectionAB sin mul add def
		XstartParallel YstartParallel){#5}
	\psline[#1](#4)(#5)
}}
%
% arrowLine[options](A)(B){n}
% #2---->---->---->---->----#3  #4-arrows inside
\def\arrowLine{\@ifnextchar[{\pst@arrowLine}{\pst@arrowLine[]}}
\def\pst@arrowLine[#1](#2)(#3)#4{{%
	\setkeys{psset}{arrowsize=4pt,arrows=->}% the defaults
	\setkeys{psset}{#1}%
	\def\pst@ThreeDplot@n{#4}
	\pst@getcoor{#2}\pst@tempa
	\pst@getcoor{#3}\pst@tempb
	\pnode(!%
	    	/YA \pst@tempa exch pop \pst@number\psyunit div def
		/XA \pst@tempa pop \pst@number\psxunit div def
		/YB \pst@tempb exch pop \pst@number\psyunit div def
		/XB \pst@tempb pop \pst@number\psxunit div def
		/dY YB YA sub \pst@ThreeDplot@n\space 1 add div def
		/dX XB XA sub \pst@ThreeDplot@n\space 1 add div def
		/Alpha dY dX atan def
		/dYOffset \psk@ThreeDplot@arrowOffset\space Alpha sin mul def
		/dXOffset \psk@ThreeDplot@arrowOffset\space Alpha cos mul def
		XA YA ){tempArrowC}
	\multido{\i=1+1}{#4}{%
		\pnode(!%
			XA dX \i\space mul add dXOffset add
			YA dY \i\space mul add dYOffset add){tempArrowB}
		\psline(tempArrowC)(tempArrowB)
		\pnode(tempArrowB){tempArrowC}
    	}
	\psline[arrows=-](tempArrowB)(#3)
}}
%
%   #1------#3------#2
\def\nodeBetween(#1)(#2)#3{%    Herbert 2003/01/05
    \pst@getcoor{#1}\pst@tempa
    \pst@getcoor{#2}\pst@tempb
    \pnode(!%
	/XA \pst@tempa pop \pst@number\psxunit div def
	/YA \pst@tempa exch pop \pst@number\psyunit div def
	/XB \pst@tempb pop \pst@number\psxunit div def
	/YB \pst@tempb exch pop \pst@number\psyunit div def
	XB XA add 2 div YB YA add 2 div){#3}
}
%
% rotateNode(A)
% (A) the node
% #2  the angle
% Herbert Voss <voss@perce.de> 2003-01-26
\def\rotateNode{\pst@rotateNode}
\def\pst@rotateNode(#1)#2{{%
    \pst@getcoor{#1}\pst@tempa
    \def\pst@ThreeDplot@angle{#2}	% the rotating angle
    \pnode(!%
    	/YA \pst@tempa exch pop \pst@number\psyunit div def
	/XA \pst@tempa pop \pst@number\psxunit div def
	YA 0 eq XA 0 eq and {0 0}{
	    /r XA dup mul YA dup mul add sqrt def
	    /AlphaOld YA XA atan def
	    /AlphaNew AlphaOld \pst@ThreeDplot@angle\space add def
	    r AlphaNew cos mul r AlphaNew sin mul
	} ifelse ){temp}
    \pnode(temp){#1}
}}
%
\def\rotateTriangle{\pst@rotateTriangle}
\def\pst@rotateTriangle(#1)(#2)(#3)#4{{%
    \rotateNode(#1){#4}
    \rotateNode(#2){#4}
    \rotateNode(#3){#4}
}}
%
\def\rotateFrame{\pst@rotateFrame}
\def\pst@rotateFrame(#1)(#2)(#3)(#4)#5{{%
    \rotateNode(#1){#5}
    \rotateNode(#2){#5}
    \rotateNode(#3){#5}
    \rotateNode(#4){#5}
}}
\catcode`\@=\PstAtCode\relax
%
%% END: pst-3dplot.tex
\endinput

%% CHANGE LOG %%
%
% 1.30	new written macros
% 1.2d	added option none for dotstyle
% 1.2c	added sperical coordinates
% 1.2b	bugfix for adding vectors with decinmal length
% 1.2a	disable the loading of pstcol.sty
